\section*{Reading map}\label{sec:reading-map}
\addcontentsline{toc}{section}{Reading map}

The manuscript is organized around a single mechanism --- the Dirac
brane stack measuring formal Hamiltonians by trace substitution into
its matrix-valued transverse coordinate --- with the rest of the
manuscript as supporting architecture.  The reading map below labels
the existing sections by their architectural role; explicit Part
numbers are non-binding bookmarks for the reader and do not enter
cross-references.

\paragraph{Part I --- Primitive physical experiment.}  The matrix
microscope statement and the seven organizing principles introduce
the experiment before any algebraic machinery.  The brane has a
matrix-valued transverse position \((\phi_1,\phi_2)\) constrained by
the conjugation moment map; substitution into a closed Hamiltonian is
the brane measurement.  See the matrix microscope section, the seven
principles (\S\ref{sec:organizing-principles}), and the introduction
\S\ref{ssec:brane-microscope-one-calculation}.

\paragraph{Part II --- Open Dirac matrix model.}  The open BV theory
of the Dirac probe stack: classical action
\(S_\partial=\int_{\R_t}\operatorname{Tr}(\phi_1\,d\phi_2+A[\phi_1,\phi_2])\),
\(GL_N\)-conjugation gauge symmetry, BV antifield
\(\psi\in\mathfrak{gl}_N[1]\) with \(Q\psi=[\phi_1,\phi_2]\), and the
Koszul derived zero-fibre model for the moment-map equation.  See
\S\ref{ssec:open-mixed-brane-states},
\S\ref{ssec:open-mixed-bft}, and the stable invariant theory
algorithm of Procesi--Razmyslov.

\paragraph{Part III --- Main local theorem.}  The matrix-microscope
measurement theorem in its two forms: \emph{Theorem A} -- algebraic
uniqueness of \(J(f)=\overline{\operatorname{Tr}f(\phi_1,\phi_2)}\)
under stable \(GL_N\)-invariance, the Dirac constraint, scalar
reduction, Darboux equivariance, and linear normalization
(Theorem~\ref{thm:Jf-coordinate-uniqueness}); \emph{Theorem B} --
the same datum as the terminal protected sector among admissible
ambient data, with the four-class universal-property obstruction
\(\operatorname{Ob}^{\mathrm{univ}}_{p,X}\)
(Theorem~\ref{thm:formal-local-universal-property}).  Proof skeleton:
stable trace generation, Dirac collapse, scalar reduction,
no-natural-correction lemma, formal completion.  See
\S\ref{ssec:local-theorem}.

\paragraph{Part IV --- Closed Hamiltonian BF sector.}  The closed BV
theory: \(\widehat{\mathfrak h}=\C[[z_1,z_2]]/\C\cdot1\), Poisson
bracket, Hamiltonian vector field convention, continuous dual
\(\mathfrak h^\vee_{\mathrm{cont}}\cong\mathcal P\) by Matlis residues,
shifted cotangent Lie algebra
\(\mathfrak g=\mathfrak h\ltimes\mathfrak h^\vee_{\mathrm{cont}}[1]\),
and the Hamiltonian BF action satisfying CME on product opens of
\(\R^2_{\mathrm{top}}\times\C^2_{\mathrm{hol}}\).  See
\S\ref{ssec:hamiltonian-bf-algebra},
\S\ref{ssec:closed-mixed-hamiltonian},
\S\ref{ssec:hamiltonian-bf-action}, and Appendix on Matlis principal
parts.

\paragraph{Part V --- Locality and factorization currents.}  Three
independent locality identities --- support \(\delta(t-t')\),
topological \(L_v=[Q,\iota_v]\), formal-holomorphic
\(L_{\partial/\partial\bar z_i}=[Q,\iota_{\partial/\partial\bar z_i}]\)
--- forcing the mixed holomorphic-topological factorization-algebra
shape on \(\R^2_{\mathrm{top}}\times\C^2_{\mathrm{hol}}\).  See
\S\ref{ssec:three-localities} and Appendix on factorization-current
conventions.

\paragraph{Part VI --- CE/PV comparison.}  The five-rung ladder of
Theorem~\ref{thm:universal-ce-pv-koszul-criterion}: cochain identity,
\(P_0\)-bracket enhancement, Maurer--Cartan kernel, bar-cobar /
Quillen, unreduced factorization-centre lift.  Type discipline
\(c_f\mapsto\theta_f\), \(u_f\mapsto O_f\); decisive computation
\(\theta_f(O_g)=O_{\{f,g\}}\); admissibility hierarchy.  See
\S\ref{ssec:reduced-ce-pv-measurement} and
\S\ref{ssec:stratified-koszul}.

\paragraph{Part VII --- Weighted-Tate analytic habitat.}  Why a
completion is necessary; weighted seminorms, weighted Hamiltonian
space \(\mathfrak h_w\), continuous dual
\(\mathfrak h^\vee_{w,\mathrm{cont}}\), convergent Casimir
\(C_{\mathfrak h,w}\), the canonical category of \(w\)-weighted
bracket-admissible RG-stable Tate completions, and the minimality of
\(B_w\) in that category.  See Tate sections T1--T5.

\paragraph{Part VIII --- Scalar anomaly and first quantum obstruction.}
The single class \(\hbar N[\bar c]\in H^2_{\mathrm{Lie}}(\bar A,\C)[[\hbar]]\)
detected on four faces (algebraic central extension, Lie cocycle,
trace-algebraic boundary, Weyl/Moyal commutator); the first quantum
example of a structural twist.  See
\S\ref{subsec:scalar-anomaly} and the scalar-anomaly transgression
theorem.

\paragraph{Part IX --- Moyal/Weyl/radial comparison.}  The
finite-certificate componentwise compatibility theorem; per-bidegree
obstruction \(\Omega^{\mathrm{rad}}_{a,b}\in\operatorname{coker}B_{a,b}\);
decorated PBW Stokes form
\(D^\square_{a,b}=C^+_{a,b}\partial_2\); reproducible finite
certificate algorithm.  See Appendix on radial parts and Moyal.

\paragraph{Part X --- QME and counterterm coherence.}  Local-functional
complex \(\mathfrak Q^\bullet_{w,\partial,\hbar}(I)\), QME
differential \(d_{\mathrm{QME}}=Q+\{I,-\}+\hbar\Delta+\cdots\); reduced
non-scalar obstruction
\([\operatorname{Ob}^{\mathrm{red}}_{w,\partial,\hbar}]\); counterterm
twist with \(\theta_3\) prototype.  See Appendix on unreduced BV/QME.

\paragraph{Part XI --- Unreduced factorization-centre lift.}  The
chain-level morphism
\(\Phi^{\mathrm{cl}}\colon\mathrm{Obs}^{\mathrm{cl}}_{\mathrm{closed},H}\big|_L
\to Z^{P_0}_{\mathrm{fact}}(\mathrm{Obs}^{\mathrm{cl}}_{\mathrm{open},N})\)
and its four-component obstruction vector
\(\operatorname{Ob}_{\mathrm{cent/QME/desc}}\); derived-centre twist.

\paragraph{Part XII --- Bochner--Martinelli current transfer and pro-Matlis replacement.}
The strict finite-window no-go: the residue-coadjoint computation
\(z_1^{N+2}\cdot\rho_{N+1,0}=-(N+2)\rho_{0,1}\) and the projected
Jacobiator \((N+1)(N+2)\rho_{0,1}\ne 0\); the obstruction
\(\operatorname{Ob}^{\Pi}_{\mathrm{BM}}\) is non-absorptive and
selects the pro-Matlis tower
\(P^\Pi=\varprojlim_N P_{\le N}\) as the coherent target replacement.

\paragraph{Part XIII --- Master deformation theory of local-to-global coherence.}
Filtered \(L_\infty\) target deformation complex
\(\mathfrak K_{\mathcal T}\), transported local Maurer--Cartan datum
\(\Theta_{\mathrm{loc}}\), curvature \(F_{\mathcal T}(\Theta_{\mathrm{loc}})\),
twisted flatness \(F_{\mathcal T}(\Theta_{\mathrm{loc}}+\tau)=0\),
three twist mechanisms (null-homotopy / structural / target
replacement), and the obstruction--twist dichotomy schema.  See
Appendix~\ref{app:master-deformation-complex}.

\paragraph{Part XIV --- Globalization from first principles.}
\v{C}ech--Weiss--Ran deformation complex on a Darboux cover,
Hamiltonian period \(\operatorname{per}_X(v)=[h_i-h_j]\), gerbe-valued
Hamiltonian descent.  Recorded as a target deformation problem in
\S\ref{ssec:master-deformation-target}; the local manuscript proves
no compact-surface global descent theorem.

\paragraph{Part XV --- Cross-volume formal-slice principle.}  The
local theorem is the formal Hamiltonian Darboux brane-probe chart of
the chain
\(\mathcal C_{\mathrm{CY}}\to E_d\text{-HolFA}\to A\to
Z^{\mathrm{der}}_{\mathrm{ch}}(A)\to\mathrm{Obs}_{\mathrm{bulk}}\),
restricted along a Dirac brane probe to the protected Hamiltonian
sector.  See cross-volume firewall
\S\ref{app:convention-contract}.

\paragraph{Part XVI --- The hidden inner principle.}  An obstruction
is the memory of structure erased by an over-strict target:
\(\hbar N[\bar c]\) remembers the constant killed by Hamiltonian
reduction; \(\mathcal P=\mathfrak h^\vee_{\mathrm{cont}}\) remembers
the distributional cotangent modes erased by polynomial PBW labels;
\(\operatorname{Ob}^{\Pi}_{\mathrm{BM}}\) remembers high transverse
modes erased by finite current truncation; \(\Omega^{\mathrm{rad}}_{a,b}\)
remembers Weyl-ordering data erased by naive trace comparison;
\(\operatorname{Ob}_{\mathrm{cent/QME/desc}}\) remembers the
homotopies erased by strict-centre language.  The closing thesis
makes the principle explicit.
